
choisirons des indéterminées différentes pour chacun des facteurs du produit de @@M0@@ par lui-même, et nous indiquerons l'indéterminée en exposant lorsque ce sera nécessaire : @@M1@@, @@M2@@, etc. Pour @@M3@@, on a donc @@M4@@. Pour définir les opérations @@M5@@ sur @@M6@@, nous allons donc nous donner un morphisme de foncteurs : @@M7@@.
Pour cela, il faut d'abord définir une filtration du foncteur @@M8@@ qui permette le recollement de 1.6.

On remarque d'abord que @@M9@@ est un idéal pour la loi @@M10@@~; pour cela on est ramené par fonctorialité au cas @@M11@@, et au produit @@M12@@, et cela résulte de 2.3.2. On prend alors pour @@M13@@ l'ensemble des séries de la forme @@M14@@, @@M15@@, @@M16@@ désignant la partie entière de @@M17@@.

L'élément définissant le morphisme @@M18@@ est @@M19@@~; le lecteur vérifiera à titre d'exercice que les conditions de 1.6 sont bien remplies (en fait, la filtration est choisie ad hoc).

Ici encore, les opérations @@M20@@ se calculent par substitution à partir de @@M21@@, et sont donc déterminées par des polynômes universels. On notera @@M22@@ le coefficient de @@M23@@ dans @@M24@@~; c'est un polynôme isobare de poids @@M25@@ par rapport aux variables @@M26@@ (vérification analogue à celle de 1.15). Il se calcule par fonctions symétriques à partir de l'élément de définition :
@@M27@@

Donnons à titre d'exemple la valeur de @@M28@@ :
